\begin{table}[t]
\tbl{\label{tab:orthogonal}}
{\begin{tcolorbox}[tabularx={p{3cm}|p{2cm}|p{6.3cm}|},enhanced,fonttitle=\bfseries\large,fontupper=\normalsize\sffamily,
colback=yellow!10!white,colframe=red!50!black,colbacktitle=blue!30!white,
coltitle=black,title=Creating an Orthogonal Space]
\hline
Command & Arguments & Purpose \\
\hline
\hline
\verb'-epsilon'  &  $\epsilon$  &  Specify the type $\epsilon$. \\
\hline
\verb'-n'  &  $n$  &  Specify the projective dimension. \\
\hline
\verb'-q'  &  $q$  &  Specify the order of the field. \\
\hline
\verb'-label_txt'  &  L  &  Set the ascii-label of the space. The label is used for things like file names etc. A default label will be used if this option is not given. \\
\hline
\verb'-label_tex'  &  L  &  Set the tex-label of the space. The label is used within latex reports. A default label will be used if this option is not given. \\
\hline
\verb'-without_group'  &    &  Do not create the orthogonal group. \\
\hline
\verb'-create_extension_fields'  &    &  Create extension fields of degree two and three of GF(q). \\
\hline
\end{tcolorbox}}
\end{table}
